\section{Process-Aware Constructibility-1}

\subsection{Scope and Central Endpoint}
Constructibility-1 is a separately versioned development release over the
six public Frontier sources. It contains 78 conditions, not 78 independent
objects. Its question is whether a proposed construction remains valid at
every prefix, exposes enough tool clearance, recovers minimally after a
blocked action, adapts to a stockout, and expresses uncertainty before
committing to a hidden world.

We retain the diagnostic vector
\begin{equation}
\mathcal{C}=(G,P,A,R,U),
\end{equation}
where $G$ is discrete geometric validity, $P$ prefix support, $A$ declared
tool accessibility, $R$ recovery, and $U$ finite-world uncertainty handling.
End-to-end success is their conjunction where applicable. We do not publish
a weighted average across task families: a compensating average would allow
strong static geometry to hide a complete recovery failure.

\begin{table}[H]
\centering\small
\begin{tabular}{lrrr}
\toprule
Family & Cases & Solver & Shortcut\\
\midrule
\tableinput{tables/constructibility-v1.tex}
\bottomrule
\end{tabular}
\caption{Constructibility-1 release audit. Solver is the number of
public-information oracle/control successes. Shortcut is always-valid,
no-operation, unavailable-part reuse, or always-commit-\texttt{h0},
respectively. The six valid sequence controls explain all six passes
of the always-valid rule.}
\label{tab:constructibility-v1}
\end{table}

\subsection{Accessibility Contract}
Every placement first satisfies integer-grid bounds, no body overlap,
at least one directly supporting stud for elevated parts, and a clear
top-down body sweep. A named west, east, north, or south approach then
requires one empty stud strip beside the part from its base upward.
The strip represents deterministic room for a release tool. Removal uses
the same condition and additionally rejects any overhead body.

This contract is intentionally stronger than body-only vertical insertion:
six constructed dead ends pass the earlier body test but fail the tool strip.
It remains a nominal proxy. It is neither a collision-checked robot trajectory
nor a calibrated statement about hand access, clutch force, or ergonomics.

For an accepted sequence of target length $T$, legal-prefix rate is
$L/T$, where $L$ is the number of steps before the first rejection.
For each accepted elevated part, support fraction is the number of directly
supported footprint studs divided by footprint area. We report the mean and
minimum over accepted prefixes. These are support-coverage diagnostics, not
force-balance stability.

\subsection{Four Task Families}
\paragraph{Sequence audit.}
Each source has one valid plan, one plan that begins with an unsupported
elevated part, and one plan that delays an enclosed part until its body still
fits but no release strip remains. The answer gives the zero-based first
failure, its typed reason, legal-prefix length, and constructible decision.
The candidate is never reordered by the evaluator.

\paragraph{Blocked recovery.}
Each source contributes three failed checkpoints. The answer diagnoses the
failed ID and clearance reason, then gives explicit removal and placement
actions. Removed instances enter inventory; only the missing failed instance
or a removed instance can be placed. The recovered checkpoint must be exact,
the number of recovery actions must match the certified minimum, and up to
twelve following target parts must then be installed legally. Thus a correct
diagnosis without repair, or a corrected checkpoint without continuation,
does not pass.

\paragraph{Stockout replanning.}
One scheduled 4-by-8 foundation plate is unavailable. The model receives a
finite stock of smaller plates and must cover exactly the unavailable
three-dimensional region at the same color, use no more than the certified
two-piece lower bound, install the replacements with declared approaches,
and execute the following construction window. Alternative seams are accepted
when they satisfy the complete contract.

\paragraph{Clarify or commit.}
Eighteen finite-world cases remain ambiguous after one geometric section
scan; eighteen matched cases add observations until one world remains.
The model returns all admissible worlds, a probability for every original
candidate under the declared uniform prior, a decision, and a next query.
Ambiguous cases require \texttt{abstain}; resolved cases require the unique
world. The unresolved next query maximizes expected entropy reduction per
unit cost, with all exact ties accepted. We report posterior $L_1$, Brier
score, false certainty, unnecessary abstention, and query correctness.
This is finite-grammar calibration, not open-world confidence estimation.

\subsection{Integrity and Reporting}
Public records omit all oracle answers and the selected hidden-world ID.
Private references, algorithm results, protocol, sources, and release files
are hash-bound. Unknown or duplicate IDs reject the submission; missing,
malformed, and truncated answers remain in each family denominator.
Raw-model, fixed-harness, and tool-assisted results must be stored and
reported separately.

All public-information controls pass. The shortcut controls in
Table~\ref{tab:constructibility-v1} fail every recovery and stockout case;
always committing to \texttt{h0} passes only 5/36 clarification cases.
These checks establish task non-degeneracy against the named shortcuts, not
learned-model discrimination.

\subsection{Unfinished Scientific Gates}
The six sources were used during task design and are not an independent test
set. No learned model, new human participant, external object, force solver,
or real robot was evaluated in this release. The next confirmation must freeze
new licensed source groups before adaptation, add matched visual process
observations, run representative models and specialized planners under equal
budgets, and obtain independent human metric review. Dynamic loading,
connection tolerances, wear, and gear/hinge function require new validated
engines or physical experiments and must not be inferred from these scores.
